\begin{thebibliography}{39}
\bibitem{tham2014}
Y.-C. Tham, X. Li, T. Y. Wong, H. A. Quigley, T. Aung, and C.-Y. Cheng, ``Global prevalence of glaucoma and projections of glaucoma burden through 2040: A systematic review and meta-analysis,'' \emph{Ophthalmology}, vol. 121, no. 11, pp. 2081--2090, 2014, doi: 10.1016/j.ophtha.2014.05.013.
\bibitem{gupta2006}
N. Gupta, L.-C. Ang, L. No\"{e}l de Tilly, L. Bidaisee, and Y. H. Y\"{u}cel, ``Human glaucoma and neural degeneration in intracranial optic nerve, lateral geniculate nucleus, and visual cortex,'' \emph{Brit. J. Ophthalmol.}, vol. 90, no. 6, pp. 674--678, 2006, doi: 10.1136/bjo.2005.086769.
\bibitem{halliday1972}
A. M. Halliday, W. I. McDonald, and J. Mushin, ``Delayed visual evoked response in optic neuritis,'' \emph{Lancet}, vol. 299, no. 7758, pp. 982--985, 1972, doi: 10.1016/S0140-6736(72)91155-5.
\bibitem{sakki2018}
H. E. A. Sakki, N. J. Dale, J. Sargent, T. Perez-Roche, and R. Bowman, ``Is there consensus in defining childhood cerebral visual impairment? A systematic review of terminology and definitions,'' \emph{Brit. J. Ophthalmol.}, vol. 102, no. 4, pp. 424--432, 2018, doi: 10.1136/bjophthalmol-2017-310694.
\bibitem{odom2016}
J. V. Odom \emph{et al.}, ``ISCEV standard for clinical visual evoked potentials: (2016 update),'' \emph{Documenta Ophthalmol.}, vol. 133, no. 1, pp. 1--9, 2016, doi: 10.1007/s10633-016-9553-y.
\bibitem{calhoun2016}
V. D. Calhoun and J. Sui, ``Multimodal fusion of brain imaging data: A key to finding the missing link(s) in complex mental illness,'' \emph{Biol. Psychiatry, Cogn. Neurosci. Neuroimaging}, vol. 1, no. 3, pp. 230--244, 2016, doi: 10.1016/j.bpsc.2015.12.005.
\bibitem{kamnitsas2017}
K. Kamnitsas \emph{et al.}, ``Efficient multi-scale 3D CNN with fully connected CRF for accurate brain lesion segmentation,'' \emph{Med. Image Anal.}, vol. 36, pp. 61--78, 2017, doi: 10.1016/j.media.2016.10.004.
\bibitem{cicek2016}
\"{O}. \c{C}i\c{c}ek, A. Abdulkadir, S. S. Lienkamp, T. Brox, and O. Ronneberger, ``3D U-Net: Learning dense volumetric segmentation from sparse annotation,'' in \emph{Proc. Med. Image Comput. Comput.-Assist. Intervent. (MICCAI)}, 2016, pp. 424--432, doi: 10.1007/978-3-319-46723-8\_49.
\bibitem{he2016}
K. He, X. Zhang, S. Ren, and J. Sun, ``Deep residual learning for image recognition,'' in \emph{Proc. IEEE Conf. Comput. Vis. Pattern Recognit. (CVPR)}, 2016, pp. 770--778, doi: 10.1109/CVPR.2016.90.
\bibitem{isensee2019}
F. Isensee \emph{et al.}, ``Automated brain extraction of multisequence MRI using artificial neural networks,'' \emph{Hum. Brain Mapp.}, vol. 40, no. 17, pp. 4952--4964, 2019, doi: 10.1002/hbm.24750.
\bibitem{hoopes2022}
A. Hoopes, J. S. Mora, A. V. Dalca, B. Fischl, and M. Hoffmann, ``SynthStrip: Skull-stripping for any brain image,'' \emph{NeuroImage}, vol. 260, Art. no. 119474, 2022, doi: 10.1016/j.neuroimage.2022.119474.
\bibitem{cardoso2022}
M. J. Cardoso \emph{et al.}, ``MONAI: An open-source framework for deep learning in healthcare,'' 2022, \emph{arXiv:2211.02701}.
\bibitem{butterworth1930}
S. Butterworth, ``On the theory of filter amplifiers,'' \emph{Exp. Wireless Wireless Eng.}, vol. 7, pp. 536--541, 1930.
\bibitem{gustafsson1996}
F. Gustafsson, ``Determining the initial states in forward-backward filtering,'' \emph{IEEE Trans. Signal Process.}, vol. 44, no. 4, pp. 988--992, 1996, doi: 10.1109/78.492552.
\bibitem{oppenheim2010}
A. V. Oppenheim and R. W. Schafer, \emph{Discrete-Time Signal Processing}, 3rd ed. Upper Saddle River, NJ, USA: Pearson, 2010.
\bibitem{baltrusaitis2019}
T. Baltru\v{s}aitis, C. Ahuja, and L.-P. Morency, ``Multimodal machine learning: A survey and taxonomy,'' \emph{IEEE Trans. Pattern Anal. Mach. Intell.}, vol. 41, no. 2, pp. 423--443, 2019, doi: 10.1109/TPAMI.2018.2798607.
\bibitem{ramachandram2017}
D. Ramachandram and G. W. Taylor, ``Deep multimodal learning: A survey on recent advances and trends,'' \emph{IEEE Signal Process. Mag.}, vol. 34, no. 6, pp. 96--108, 2017, doi: 10.1109/MSP.2017.2738401.
\bibitem{vaswani2017}
A. Vaswani \emph{et al.}, ``Attention is all you need,'' in \emph{Proc. Adv. Neural Inf. Process. Syst. (NeurIPS)}, vol. 30, 2017, pp. 5998--6008.
\bibitem{dosovitskiy2021}
A. Dosovitskiy \emph{et al.}, ``An image is worth 16\texttimes16 words: Transformers for image recognition at scale,'' in \emph{Proc. Int. Conf. Learn. Represent. (ICLR)}, 2021.
\bibitem{tsai2019}
Y.-H. H. Tsai, S. Bai, P. P. Liang, J. Z. Kolter, L.-P. Morency, and R. Salakhutdinov, ``Multimodal transformer for unaligned multimodal language sequences,'' in \emph{Proc. 57th Annu. Meeting Assoc. Comput. Linguistics}, 2019, pp. 6558--6569, doi: 10.18653/v1/P19-1656.
\bibitem{lu2019}
J. Lu, D. Batra, D. Parikh, and S. Lee, ``ViLBERT: Pretraining task-agnostic visiolinguistic representations for vision-and-language tasks,'' in \emph{Proc. Adv. Neural Inf. Process. Syst. (NeurIPS)}, vol. 32, 2019, pp. 13--23.
\bibitem{selvaraju2017}
R. R. Selvaraju, M. Cogswell, A. Das, R. Vedantam, D. Parikh, and D. Batra, ``Grad-CAM: Visual explanations from deep networks via gradient-based localization,'' in \emph{Proc. IEEE Int. Conf. Comput. Vis. (ICCV)}, 2017, pp. 618--626, doi: 10.1109/ICCV.2017.74.
\bibitem{jain2019}
S. Jain and B. C. Wallace, ``Attention is not explanation,'' in \emph{Proc. Conf. North Amer. Chapter Assoc. Comput. Linguistics: Human Lang. Technol. (NAACL-HLT)}, 2019, pp. 3543--3556, doi: 10.18653/v1/N19-1357.
\bibitem{wiegreffe2019}
S. Wiegreffe and Y. Pinter, ``Attention is not not explanation,'' in \emph{Proc. Conf. Empirical Methods Natural Lang. Process. 9th Int. Joint Conf. Natural Lang. Process. (EMNLP-IJCNLP)}, 2019, pp. 11--20, doi: 10.18653/v1/D19-1002.
\bibitem{adebayo2018}
J. Adebayo, J. Gilmer, M. Muelly, I. Goodfellow, M. Hardt, and B. Kim, ``Sanity checks for saliency maps,'' in \emph{Proc. Adv. Neural Inf. Process. Syst. (NeurIPS)}, vol. 31, 2018, pp. 9505--9515.
\bibitem{vanessen2013}
D. C. Van Essen, S. M. Smith, D. M. Barch, T. E. J. Behrens, E. Yacoub, and K. Ugurbil, ``The WU-Minn Human Connectome Project: An overview,'' \emph{NeuroImage}, vol. 80, pp. 62--79, 2013, doi: 10.1016/j.neuroimage.2013.05.041.
\bibitem{benson2018}
N. C. Benson \emph{et al.}, ``The Human Connectome Project 7 Tesla retinotopy dataset: Description and population receptive field analysis,'' \emph{J. Vis.}, vol. 18, no. 13, Art. no. 23, 2018, doi: 10.1167/18.13.23.
\bibitem{lamontagne2019}
P. J. LaMontagne \emph{et al.}, ``OASIS-3: Longitudinal neuroimaging, clinical, and cognitive dataset for normal aging and Alzheimer disease,'' \emph{medRxiv}, 2019, doi: 10.1101/2019.12.13.19014902.
\bibitem{gramfort2013}
A. Gramfort \emph{et al.}, ``MEG and EEG data analysis with MNE-Python,'' \emph{Front. Neurosci.}, vol. 7, Art. no. 267, 2013, doi: 10.3389/fnins.2013.00267.
\bibitem{goldberger2000}
A. L. Goldberger \emph{et al.}, ``PhysioBank, PhysioToolkit, and PhysioNet: Components of a new research resource for complex physiologic signals,'' \emph{Circulation}, vol. 101, no. 23, pp. e215--e220, 2000, doi: 10.1161/01.CIR.101.23.e215.
\bibitem{paszke2019}
A. Paszke \emph{et al.}, ``PyTorch: An imperative style, high-performance deep learning library,'' in \emph{Proc. Adv. Neural Inf. Process. Syst. (NeurIPS)}, vol. 32, 2019, pp. 8024--8035.
\bibitem{virtanen2020}
P. Virtanen \emph{et al.}, ``SciPy 1.0: Fundamental algorithms for scientific computing in Python,'' \emph{Nature Methods}, vol. 17, no. 3, pp. 261--272, 2020, doi: 10.1038/s41592-019-0686-2.
\bibitem{harris2020}
C. R. Harris \emph{et al.}, ``Array programming with NumPy,'' \emph{Nature}, vol. 585, no. 7825, pp. 357--362, 2020, doi: 10.1038/s41586-020-2649-2.
\bibitem{hunter2007}
J. D. Hunter, ``Matplotlib: A 2D graphics environment,'' \emph{Comput. Sci. Eng.}, vol. 9, no. 3, pp. 90--95, 2007, doi: 10.1109/MCSE.2007.55.
\bibitem{zhou2024camf}
Z. Zhou, A. Orlichenko, G. Qu, Z. Fu, V. D. Calhoun, Z. Ding, and Y.-P. Wang, ``An interpretable cross-attentive multi-modal MRI fusion framework for schizophrenia diagnosis,'' 2024, \emph{arXiv:2404.00144}.
\bibitem{dhar2025drifa}
J. Dhar, N. Zaidi, M. Haghighat, S. Roy, P. Goyal, A. Alavi, and V. Kumar, ``Multimodal fusion learning with dual attention for medical imaging,'' in \emph{Proc. IEEE/CVF Winter Conf. Appl. Comput. Vis. (WACV)}, 2025, pp. 4362--4371, doi: 10.1109/WACV61041.2025.00428.
\bibitem{hu2018}
J. Hu, L. Shen, and G. Sun, ``Squeeze-and-excitation networks,'' in \emph{Proc. IEEE Conf. Comput. Vis. Pattern Recognit. (CVPR)}, 2018, pp. 7132--7141, doi: 10.1109/CVPR.2018.00745.
\bibitem{liuzzi2022ds003568}
L. Liuzzi, K. Chang, H. Keren, C. Zheng, D. Saha, D. Nielson, and A. Stringaris, ``Mood induction in MDD and healthy adolescents,'' OpenNeuro dataset ds003568, v1.0.4, 2023, doi: 10.18112/openneuro.ds003568.v1.0.4 (see also L. Liuzzi \emph{et al.}, ``Electrophysiological correlates of mood and reward dynamics in human adolescents,'' \emph{bioRxiv}, 2021, doi: 10.1101/2021.03.04.433969).
\bibitem{ds007181}
L. Li, Q. Han, H. Bai, X. Zhang, Y. Liu, and C. He, ``Structural MRI, resting-state fMRI, and PSG/EEG dataset of zoster-associated neuralgia,'' OpenNeuro dataset ds007181, v1.0.1, 2026, doi: 10.18112/openneuro.ds007181.v1.0.1.
\end{thebibliography}
